An applied action records what an agent system did, but does not by itself establish that the model chose it, could have done otherwise, or influenced another agent. We study these distinctions in twenty complete MonopolyBench games containing 6,989 decisions and 8,679 response attempts. Our procedure links each claim to its historical state, response provenance, delivered information and a specified endpoint. For feasibility claims, unchanged engine execution supplies legal witnesses or scoped impossibility certificates. A portable audit covers all sixty eliminations and sixty-six engine defaults: fifty-six eliminations cannot pay even under a generous asset bound, three admit legal payment, and one is blocked by hotel liquidation. At least one recorded construction target is reachable in sixty-three defaulted decisions after atomic expansion and reordering, although every well-formed submitted target is invalid as encoded. Across 663 END actions where the construction tool was listed, sixty-six are defaults and twenty-four have detected content/tool name conflicts. Neither class supports an unqualified investment-restraint label. Delivery reconstruction further distinguishes canonical speech from recipient information. Selected negotiations show explicit correction with acceptance and textual echo without identified uptake of a false minimum. The cohort provides executable counterexamples to particular behavioral inferences, rather than a model ranking or deception rate. All 50,918 original events replay, and the audit reproduces in a clean environment. This establishes a bounded method for checking claims about recorded trajectories; causal persuasion and capability--alignment trends require additional experiments.
